% 依据英文版表1逐项翻译，数值与条件沿用原稿。
\begingroup
\ReportTableSetup
\begin{longtable}{|L{0.15\linewidth}|L{0.23\linewidth}|L{\dimexpr0.62\linewidth-6\tabcolsep-4\arrayrulewidth\relax}|}
\caption{各州数据中心税收优惠状态}\label{tab:state-incentives}\\
\hline
\textbf{州} & \textbf{当前状态} & \textbf{主要变化} \\ \hline
\endfirsthead
\caption[]{各州数据中心税收优惠状态（续）}\\
\hline
\textbf{州} & \textbf{当前状态} & \textbf{主要变化} \\ \hline
\endhead
亚利桑那州 & 暂停新申请 & 自2026年7月1日至2029年6月30日，亚利桑那州商务管理局不得受理新计算机数据中心的申请。 \\ \hline
俄亥俄州 & 暂停新申请 & 俄亥俄州税收抵免管理局已暂停审议新的数据中心税收豁免请求。 \\ \hline
伊利诺伊州 & 暂停新协议 & 自2026年7月1日起，该州停止处理数据中心激励计划的新申请。 \\ \hline
马萨诸塞州 & 暂停新申请 & 自2026年6月25日起，该州暂停受理合资格数据中心销售和使用税豁免申请。 \\ \hline
华盛顿州 & 缩小免税范围 & 自2026年7月1日起，服务器更换设备不再适用原有豁免；不再签发新的农村地区翻新证书，城市地区翻新项目也不再符合资格。 \\ \hline
弗吉尼亚州 & 仍可免税；新增用电税 & 销售和使用税豁免继续有效；自2026年7月1日至2028年6月30日，按每千瓦时0.011美元征收数据中心用电税。 \\ \hline
路易斯安那州 & 仍可免税；条件待定 & 2026年6月的一项行政命令要求，为数据中心销售税优惠建立电费缴纳者和社区保护框架，具体标准由相关主管部门制定中。 \\ \hline
得克萨斯州 & 仍可免税；政策变化风险较高 & 该州要求全面审查数据中心电网互联项目，并提出与立法机关合作，取消过时或不必要的数据中心优惠。 \\ \hline

\end{longtable}
\endgroup
